@@M0@@ of @@M1@@; an element @@M2@@ of @@M3@@ determines a morphism of functors if and only if it satisfies:
@@M4@@

The reader will readily verify that, apart from complications of notation, it suffices to copy the proof of 1.6.

\sgaStableTarget{sga6.e05.para.1-8}1.8. We denote by @@M5@@ the category of @@M6@@-algebras, and by
@@M7@@
the functor that assigns to a @@M8@@-algebra @@M9@@ the @@M10@@-algebra @@M11@@ endowed with the coarse filtration; we denote by
@@M12@@
the functor ``forget the filtration.'' These functors satisfy the following relations:
@@M13@@
(cf. 1.3, ex. 2)

\sgaStableTarget{sga6.e05.lem.1-9}\textbf{Lemma 1.9.} \uline{Let} @@M14@@ \uline{be a functor on} @@M15@@; \uline{then the canonical homomorphism}:
@@M16@@
